\subsection{紧空间在连续映射下的象}

定理 2. 若 \(f\) 是拟紧空间 \(\mathbf{X}\) 到拓扑空间 \(\mathbf{X}'\) 中的一个连续映射，则集合 \(f(\mathbf{X})\) 是拟紧的。

设 \(\Re\) 是 \(f(\mathbf{X})\) 的一个由 \(\mathbf{X}'\) 中开集组成的覆盖；则 \(\overline{f}^{1}(\Re)\) 是 \(\mathbf{X}\) 的一个开覆盖（\(\S 2\)，第 1 目，定理 1）；于是存在有限子集 \(\mathfrak{S}\)（\(\Re\) 的），使 \((\mathfrak{S}\overline{f}^{1})\) 是 \(\mathbf{X}\) 的一个覆盖；但这时 \(\mathfrak{S}\) 就是 \(f(\mathbf{X})\) 的一个覆盖，定理得证。

推论 1. 设 \(f\) 是拓扑空间 \(\mathbf{X}\) 到豪斯多夫空间 \(\mathbf{X}'\) 中的一个连续映射。则在 \(f\) 之下，\(\mathbf{X}\) 中任一拟紧（相应地，相对拟紧）集的象是 \(\mathbf{X}'\) 中的一个紧（相应地，相对紧）集。

推论 2. 每个连续映射 \(f\)（拟紧空间 \(X\) 到豪斯多夫空间 \(X'\) 中的）都是闭映射。若此外 \(f\) 还是双射，则 \(f\) 是一个同胚。

这由第 3 目的推论 1 与命题 4 立即得出。

特别地：

推论 3. 比拟紧空间的拓扑更粗的豪斯多夫拓扑必与后者一致。

推论 4. 设 X 是一个拓扑空间，R 是 X 上的一个豪斯多夫等价关系。

\begin{enumerate}
\def\labelenumi{\alph{enumi})}
\item
  若 X 中存在一个与每个模 R 等价类都相交的拟紧集 K，则 X/R 是紧的，且 \(K/R_{K}\) 到 X/R 上的典范映射是一个同胚。
\item
  若 K 与每个等价类还只交于一点，则 K 是 X 关于关系 R 的一个连续截面（§ 3，第 5 目）。
\end{enumerate}

设 \(f\) 是限制到 \(K\) 上的典范映射 \(X \to X / R\)。既然 \(X / R\) 是豪斯多夫的，由推论 1 得出 \(X / R\) 是紧的；又由推论 2 得出 \(f\) 是闭映射；于是双射 \(K / R_K \to X / R\)（伴随于 \(f\)）是一个同胚（§ 5，第 2 目，命题 3）。这就处理了 a)；b) 立即得出，因为这时有 \(K / R_K = K\)。

